% ID: SEL-D0020-OA-E04
% ESTADO: cotejada. Segunda pasada visual completa y borrador compilado revisados.
% Fuente: Selectividad/matematicas_aplicadas/reserva b/Examen-Matemáticas Aplicadas a las Ciencias Sociales.pdf
% SHA-256: 37584a9604930b74008a6bda88ef3d61fa926d5e6a540ab9013569258f772256
% Página: 1; opción A; ejercicio original 4.

El gasto que tienen los jóvenes durante un fin de semana es una variable aleatoria que sigue una distribución Normal de media $\mu$ desconocida y desviación típica igual a 6 euros.
\begin{enumerate}[label=\alph*)]
\item \textbf{(1.5 puntos)} Se toma una muestra aleatoria simple y se obtiene que el intervalo de confianza al 95\% para la media $\mu$ es $(24.47,26.43)$. Calcule el valor de la media muestral y el tamaño de la muestra elegida.
\item \textbf{(1 punto)} Escogida otra muestra de tamaño 49 para estimar $\mu$, calcule el error máximo cometido para esa estimación con un nivel de confianza del 97\%.
\end{enumerate}
